\begin{figure}[H]
    \centering
\begin{tikzpicture}[
    scale=1.5,
    cota/.style={<->, >=stealth, thin, gray!80!black},  % Estilo das linhas de cota
    retangulo/.style={draw=gray!80!black, thick},        % Estilo do retângulo
    eixo/.style={->, >=stealth, thick, gray!80!black}   % Estilo dos eixos x e y
]

    % 1. Coordenadas do retângulo
    \coordinate (A) at (0, 0);
    \coordinate (B) at (2, 0);
    \coordinate (C) at (2, 3);
    \coordinate (D) at (0, 3);

    % 2. Hachura no lado esquerdo (fora do retângulo, sem borda)
    \path[pattern={Lines[angle=45, distance=6pt, line width=0.5pt]}, pattern color=gray!60] (2, 0) rectangle (2.3, 3);

    % 3. Desenho do Retângulo
    \draw[retangulo] (A) rectangle (C);

    % Linha branca tracejada sobre a parede esquerda (A -- D)
    \draw[dashed, white] (A) -- (D);

    % 4. Eixos de Coordenadas x e y (Origem no canto inferior esquerdo, dentro do retângulo)
    \draw[eixo] (0.2, 0.2) -- (0.8, 0.2) node[above, font=\small\bfseries] {$x$};
    \draw[eixo] (0.2, 0.2) -- (0.2, 0.8) node[right, font=\small\bfseries] {$y$};

    % --- COTA HORIZONTAL (Texto abaixo da linha) ---
    \def\offsetH{0.4}
    \draw[cota] (0, -\offsetH) -- (2, -\offsetH) 
        node[midway, below=3pt, font=\small] {$153\text{ cm}$};

    % --- COTA VERTICAL (Afastada a 0.6 para não sobrepor a hachura) ---
    \def\offsetV{0.6}
    \draw[cota] (-\offsetV, 0) -- (-\offsetV, 3) 
        node[midway, left=3pt, font=\small, rotate=90, anchor=south] {$556 \text{ cm}$};

\end{tikzpicture}
    %\caption{Caption}
    %\label{fig:placeholder}
\end{figure}